\clearpage
\begin{examplebox}

\textbf{\textcolor{CExample}{例 1｜最低门槛：已知一个端点，秒求另一个}}\par\smallskip
\textbf{题目：}抛物线 $y^2=4x$ 的一条焦点弦为 $AB$，若 $A$ 的纵坐标为 $-2$，求 $B$ 的坐标。\par\smallskip
\textbf{抓信号：}“焦点弦 + 已知一端纵坐标”$\Rightarrow$ 先用纵积，后用抛物线方程。
\[
 p=2,\quad y_Ay_B=-4
 \ \Longrightarrow\ (-2)y_B=-4
 \ \Longrightarrow\ y_B=2.
\]
\[
 x_B=\frac{y_B^2}{4}=1
 \quad\Longrightarrow\quad \boxed{B(1,2)}.
\]
\textbf{不必}设直线、不必联立。

\par\smallskip\textbf{回验：}$A(1,-2)$、$B(1,2)$ 的连线 $x=1$ 正是通径，横积 $1$、纵积 $-4$。\medskip\hrule\medskip
\textbf{\textcolor{CExample}{例 2｜迁移：焦半径转成横坐标}}\par\smallskip
\textbf{题目：}抛物线 $y^2=8x$ 的焦点弦 $AB$ 满足 $|AF|=5$，求 $|BF|$。\par\smallskip
\textbf{入手链：}已知焦半径 $\to$ 端点横坐标 $\to$ 横积 $\to$ 另一焦半径。
\[
 p=4,\qquad |AF|=x_A+2=5
 \quad\Longrightarrow\quad x_A=3.
\]
\[
 x_Ax_B=\frac{p^2}{4}=4
 \quad\Longrightarrow\quad x_B=\frac43.
\]
\[
 \boxed{|BF|=x_B+2=\frac{10}{3}}.
\]
\textbf{核对：}$\frac1{5}+\frac1{10/3}=\frac12=\frac2p$，与焦点弦倒数性质相符。


\end{examplebox}
\clearpage
\begin{examplebox}

\textbf{\textcolor{CExample}{例 3｜应用：两端点的切线互相垂直}}\par\smallskip
\textbf{题目：}在 $y^2=4x$ 上，过焦点 $F(1,0)$ 的直线
$l:x=my+1$ 交抛物线于 $A,B$。求证两点处切线互相垂直，并求其交点 $T$。\par\smallskip
\textbf{第一步：}由 $2p=4$ 知 $p=2$，直接用焦点弦结论 $y_1y_2=-4$。
两点处切线斜率
\[
 k_1=\frac2{y_1},\quad k_2=\frac2{y_2}
 \quad\Longrightarrow\quad
 \boxed{k_1k_2=\frac4{y_1y_2}=-1}.
\]
\textbf{第二步：}由 $x=my+1$ 代入 $y^2=4x$，得
$y^2-4my-4=0$，所以 $y_1+y_2=4m$。
\[
 \boxed{T\!\left(-1,\frac{y_1+y_2}{2}\right)=(-1,2m)}.
\]
\textbf{结论：}两切线垂直，交点始终落在准线 $x=-1$ 上。


\end{examplebox}
\clearpage
\begin{examplebox}
\begin{center}\begin{tikzpicture}[font=\small,>=stealth]
\draw[rounded corners=3pt,gray!45] (0,0) rectangle (14.2,12.5);

\node at (7.1,11.95) {切线推广：两切线成直角，交点恰好落在准线上};
\begin{scope}[shift={(3.1,4.2)},scale=1.05]
\draw[->,gray] (-1.65,0)--(4.9,0) node[above] {$x$};
\draw[->,gray] (0,-2.2)--(0,5.15) node[left] {$y$};
\draw[violet,dashed] (-1,-2.2)--(-1,5.05);
\node[violet,rotate=90] at (-1.5,3.8) {准线 $x=-1$};
\draw[teal!75!black,thick,samples=120,domain=-2.2:4.5,variable=\t] plot ({\t*\t/4},{\t});
\draw[gray!75,thick] (4,4)--(.25,-1);
\draw[blue!75!black,very thick] (-1,1.5)--(4,4);
\draw[orange!85!black,very thick] (-1,1.5)--(.25,-1);
\draw[black] (-.68695,1.65652)--(-.53043,1.34348)--(-.84348,1.18695);
\fill (-1,1.5) circle(2pt) node[left=8pt] {$T$};
\fill (4,4) circle(2pt) node[above left=7pt] {$A$};
\fill (.25,-1) circle(2pt) node[below left=7pt] {$B$};
\fill (1,0) circle(2pt) node[below right=6pt] {$F$};
\end{scope}
\node[anchor=west,align=left] at (9.1,7.15)
 {$A(4,4)$\\[12pt]$B(1/4,-1)$\\[12pt]$T(-1,3/2)$\\[16pt]
 蓝色切线：$k_A=1/2$\\[12pt]橙色切线：$k_B=-2$\\[14pt]$k_Ak_B=-1$};
\node at (7.1,1.4) {灰线是焦点弦 $AB$；蓝、橙两线是端点处的切线。};
\node at (7.1,.55) {$y^2=4x,\quad p=2,\quad y_Ay_B=-4\quad\Longrightarrow\quad k_Ak_B=-1$};
\end{tikzpicture}
\end{center}

\textbf{\textcolor{CExample}{例 4｜反向应用：只给一端，不给焦点弦方程}}\par\smallskip
\textbf{题目：}抛物线 $y^2=4x$ 上有点 $A(4,4)$，过焦点 $F(1,0)$ 作直线与抛物线再交于 $B$，求 $B$ 以及两切线交点 $T$。\par\smallskip
\[
 p=2,\qquad y_B=-\frac{p^2}{y_A}=-1,\qquad
 x_B=\frac{(-1)^2}{4}=\frac14.
\]
\[
 \boxed{B\!\left(\frac14,-1\right)},\qquad
 \boxed{T\!\left(-1,\frac{4+(-1)}{2}\right)=\left(-1,\frac32\right)}.
\]
\textbf{检验三点共线：}$\frac{4-0}{4-1}=\frac43$，
$\frac{-1-0}{\frac14-1}=\frac43$，成立。

\par\smallskip\textbf{回验切线：}$4y=2(x+4)$ 与 $-y=2(x+1/4)$ 都通过 $T(-1,3/2)$，斜率分别为 $1/2,-2$。
\end{examplebox}
\clearpage
\begin{examplebox}

\textbf{\textcolor{CExample}{高价值变式｜换成上开口，还能直接照抄吗？}}\par\smallskip
\textbf{题目：}在 $x^2=6y$ 中，过焦点的弦端点为 $A(x_1,y_1),B(x_2,y_2)$，求两个坐标积。\par\smallskip
\textbf{答案：}此时 $2p=6\Rightarrow p=3$，\textbf{交换横纵角色}：
\[
 \boxed{x_1x_2=-9},\qquad
 \boxed{y_1y_2=\frac94}.
\]
\textbf{迁移点：}不是见到“焦点弦”就机械写 $x_1x_2=\frac{p^2}{4}$；必须先看\textbf{哪个变量被平方}。


\end{examplebox}
